\chapter{SABR and Smile Dynamics}\label{ch:dv:sabr-and-smile-dynamics}

In 2002 four authors working on an interest-rate desk published a short paper with an
uncomfortable finding. The model many desks used to fit swaption smiles, \omterm{def:dv:local-volatility:model}{local volatility},
predicted that when the forward rate rose the smile would move the other way, towards lower
strikes; the market's smile moved with the forward. A model that gets the direction wrong
gets the delta wrong, and the authors showed that its hedges could be worse than hedging with
no smile at all. Their replacement had four parameters, a closed-form approximation for the
implied volatility that a spreadsheet could evaluate, and the right dynamics. SABR became the
market standard for interest-rate options and a common one for currency options. This chapter
states the model and its formula, explains the \omterm{def:dv:sabr-and-smile-dynamics:backbone}{backbone} that governs its dynamics, derives the
delta that its dynamics imply, measures by simulation what that delta is worth, and shows where
the formula breaks.

\section{The SABR model}

\begin{definition}[SABR model]\label{def:dv:sabr-and-smile-dynamics:sabr}
The \emph{SABR model}\index{SABR model} (stochastic alpha, beta, rho) describes a forward $F_t$,
driftless under its forward measure, and its volatility $\alpha_t$:
\[
dF_t=\alpha_tF_t^{\beta}\,dW^1_t,\qquad d\alpha_t=\nu\,\alpha_t\,dW^2_t,\qquad
d\langle W^1,W^2\rangle_t=\rho\,dt,
\]
with an exponent $\beta\in[0,1]$, a correlation $\rho$ and a \omterm{def:dv:stochastic-volatility:sv}{volatility of volatility} $\nu$; the
initial value $\alpha=\alpha_0$ is the fourth parameter.
\end{definition}

It is a \omterm{def:dv:stochastic-volatility:sv}{stochastic volatility model} in the sense of chapter 10, with two differences from Heston:
the volatility is lognormal and does not mean-revert, which suits a model calibrated one expiry at a
time, and the forward has a constant-elasticity diffusion $F^\beta$, which interpolates between a
normal model ($\beta=0$, the Bachelier model of One Quant Book 2, chapter 13) and a lognormal one
($\beta=1$). Interest-rate desks calibrate a separate SABR to each expiry and tenor of the volatility
cube; the normal and shifted versions they use when rates approach zero are One Quant Book 6,
chapter 5.

\section{The asymptotic formula}

SABR has no closed-form price. Its authors derived, by singular perturbation in the time to expiry
and the \omterm{def:dv:implied-volatility-and-its-surface:surface}{log-moneyness}, an explicit approximation of the Black implied volatility.

\begin{definition}[Hagan formula]\label{def:dv:sabr-and-smile-dynamics:hagan}
The \emph{Hagan formula}\index{Hagan formula} for the SABR Black volatility at strike $K$ and expiry
$T$ is
\[
\sigma_B(K)=\frac{\alpha}{(FK)^{\frac{1-\beta}2}\bigl(1+\frac{(1-\beta)^2}{24}\ell^2+
\frac{(1-\beta)^4}{1920}\ell^4\bigr)}\,\frac{z}{x(z)}\,\bigl(1+cT\bigr),
\]
\[
c=\frac{(1-\beta)^2\alpha^2}{24(FK)^{1-\beta}}+\frac{\rho\beta\nu\alpha}{4(FK)^{\frac{1-\beta}2}}
+\frac{2-3\rho^2}{24}\nu^2,
\]
with $\ell=\ln(F/K)$, $z=\frac\nu\alpha(FK)^{\frac{1-\beta}2}\ell$ and
$x(z)=\ln\bigl((\sqrt{1-2\rho z+z^2}+z-\rho)/(1-\rho)\bigr)$.
\end{definition}

At the money $z/x(z)\to1$ and the formula reduces to
$\sigma_{\mathrm{ATM}}=\alpha F^{\beta-1}\bigl(1+(\dots)T\bigr)$, a cubic in $\alpha$. Desks therefore
parametrise SABR by the at-the-money volatility they observe, solve the cubic for $\alpha$, fix $\beta$
by convention or from the \omterm{def:dv:sabr-and-smile-dynamics:backbone}{backbone} (next section), and fit $\rho$ and $\nu$ to the smile. The normal
(Bachelier) volatility the rates market quotes is the same price read with the normal model.

\begin{example}[A three-percent forward]\label{ex:dv:sabr-and-smile-dynamics:rates}
A one-year option on a forward rate of 3\%, at-the-money Black volatility 20\%, $\beta=0.5$, $\rho=-0.6$,
$\nu=0.5$: the cubic gives $\alpha=0.0346$. The Black smile runs from 28.8\% at a 2\% strike to 15.8\% at
4\%; the same prices read as normal volatilities run from 70.8 to 55.0 basis points a year, 59.9 at the money. The
normal smile is far less skewed: much of the Black skew is the lognormal model's own scaling.
\end{example}

\begin{example}[SABR on an equity smile]\label{ex:dv:sabr-and-smile-dynamics:equity}
With $\beta=1$ (a lognormal \omterm{def:dv:sabr-and-smile-dynamics:backbone}{backbone}), SABR fitted to the one-year smile of chapter 9's surface returns
$\alpha=0.194$, $\rho=-0.68$, $\nu=0.72$, with a root-mean-square error of 0.07 volatility point and at most
0.15 point, at the 70 strike (\cref{fig:dv:sabr-and-smile-dynamics:equity}). One expiry, three free
parameters, a good fit: SABR is a smile interpolator as much as a model.
\end{example}

\begin{omfigure}
\begin{tikzpicture}
\begin{axis}[width=10.5cm,height=4.8cm,
  xlabel={strike},ylabel={1-year implied volatility (\%)},
  tick label style={font=\small},label style={font=\small},
  xmin=68,xmax=132,ymin=14,ymax=30,grid=major,grid style={gray!25},
  legend columns=2,legend style={font=\footnotesize,at={(0.5,-0.3)},anchor=north,draw=none,fill=none,/tikz/every even column/.append style={column sep=8pt}},
  table/col sep=comma]
\addplot[omDef,only marks,mark=*,mark size=1.8pt] table[x=strike,y=market]{figdata/derivatives/11-sabr-and-smile-dynamics/equity_fit.csv};
\addplot[omThm,very thick] table[x=strike,y=sabr]{figdata/derivatives/11-sabr-and-smile-dynamics/equity_fit.csv};
\legend{market (chapter 9's surface),SABR with $\beta=1$}
\end{axis}
\end{tikzpicture}

\omcaption{SABR with a lognormal \omterm{def:dv:sabr-and-smile-dynamics:backbone}{backbone} fitted to one equity expiry: 13 strikes, errors of at most 0.15
volatility point. Data: the tutorial.}\label{fig:dv:sabr-and-smile-dynamics:equity}
\end{omfigure}

\section{Backbone and smile dynamics}

\begin{definition}[Backbone]\label{def:dv:sabr-and-smile-dynamics:backbone}
The \emph{backbone}\index{backbone} of a smile model is the curve traced by the at-the-money implied
volatility as the forward moves, all other state variables fixed. In SABR it is
$\sigma_{\mathrm{ATM}}(F)\approx\alpha F^{\beta-1}$: flat for $\beta=1$, falling like $1/F$ for $\beta=0$.
\end{definition}

The \omterm{def:dv:sabr-and-smile-dynamics:backbone}{backbone} is how SABR encodes dynamics. When the forward moves, the whole smile moves with it: the
at-the-money point slides along the \omterm{def:dv:sabr-and-smile-dynamics:backbone}{backbone} and the smile's shape travels with the new forward
(\cref{fig:dv:sabr-and-smile-dynamics:backbone}). That is the behaviour Hagan and his coauthors observed in
rates markets and the opposite of \omterm{def:dv:local-volatility:model}{local volatility} (chapter 9), whose smile moves against the forward. Two
parameters produce skew: $\beta<1$ gives a skew along the \omterm{def:dv:sabr-and-smile-dynamics:backbone}{backbone}, $\rho<0$ a skew relative to it. They
are nearly interchangeable for fitting one smile, and distinguishable only by dynamics: $\beta$ is chosen from
how at-the-money volatility has moved with the forward historically, or by convention.

\begin{omfigure}
\begin{tikzpicture}
\begin{groupplot}[group style={group size=2 by 1,horizontal sep=1.5cm},
  width=6.6cm,height=4.8cm,
  tick label style={font=\small},label style={font=\small},grid=major,grid style={gray!25},table/col sep=comma]
\nextgroupplot[xlabel={forward (\%)},ylabel={at-the-money volatility (\%)},xmin=2,xmax=4,ymin=14,ymax=31,
  title={\small backbones},
  legend columns=3,legend style={font=\footnotesize,at={(0.5,-0.32)},anchor=north,draw=none,fill=none}]
\addplot[omThm,very thick] table[x=fwd,y=b0]{figdata/derivatives/11-sabr-and-smile-dynamics/backbone.csv};
\addplot[omDef,thick] table[x=fwd,y=b05]{figdata/derivatives/11-sabr-and-smile-dynamics/backbone.csv};
\addplot[omProp,thick,dashed] table[x=fwd,y=b1]{figdata/derivatives/11-sabr-and-smile-dynamics/backbone.csv};
\legend{$\beta=0$,$\beta=0.5$,$\beta=1$}
\nextgroupplot[xlabel={strike (\%)},ylabel={Black volatility (\%)},xmin=1.5,xmax=5,ymin=14,ymax=36,
  title={\small smiles as the forward moves ($\beta=0.5$)},
  legend columns=3,legend style={font=\footnotesize,at={(0.5,-0.32)},anchor=north,draw=none,fill=none}]
\addplot[omThm,thick] table[x=k,y=f025]{figdata/derivatives/11-sabr-and-smile-dynamics/smiles.csv};
\addplot[omDef,very thick] table[x=k,y=f03]{figdata/derivatives/11-sabr-and-smile-dynamics/smiles.csv};
\addplot[omProp,thick,dashed] table[x=k,y=f035]{figdata/derivatives/11-sabr-and-smile-dynamics/smiles.csv};
\legend{$F=2.5\%$,$F=3\%$,$F=3.5\%$}
\end{groupplot}
\end{tikzpicture}

\omcaption{SABR dynamics. Left: the at-the-money volatility as the forward moves, for three values of
$\beta$ with $\alpha$ set to give 20\% at a 3\% forward. Right: the smile at three forwards ($\beta=0.5$,
$\rho=-0.6$, $\nu=0.5$): it moves with the forward, its minimum travelling to the right as the forward rises.
Data: the chapter's code.}\label{fig:dv:sabr-and-smile-dynamics:backbone}
\end{omfigure}

\section{Hedging under a smile model}

A smile model's delta is the change of the option's value when the underlying moves, with the model's own
view of what the other state variables do meanwhile. Three answers are in use.

\begin{itemize}
\item The \emph{Black delta} freezes the option's implied volatility: $\Delta_{\mathrm{BS}}$.
\item \emph{Hagan's delta} moves the volatility along the smile as $F$ moves, with $\alpha$ fixed:
$\Delta_{\mathrm{BS}}+\mathcal V\,\partial\sigma_B/\partial F$.
\item The third takes into account that $\alpha$ itself tends to move when $F$ moves, since the two
Brownian motions are correlated.
\end{itemize}

\begin{definition}[Minimum-variance delta]\label{def:dv:sabr-and-smile-dynamics:mvdelta}
The \emph{minimum-variance delta}\index{minimum-variance delta} of an option in a \omterm{def:dv:stochastic-volatility:sv}{stochastic volatility model}
is the position in the underlying that minimises the variance of the hedged position's instantaneous P\&L.
In SABR (Bartlett's delta) it is
\[
\Delta_{\mathrm{MV}}=\Delta_{\mathrm{BS}}+\mathcal V\Bigl(\frac{\partial\sigma_B}{\partial F}
+\frac{\partial\sigma_B}{\partial\alpha}\,\frac{\rho\nu}{F^{\beta}}\Bigr).
\]
\end{definition}

\begin{proposition}[Why that delta]\label{prop:dv:sabr-and-smile-dynamics:mv}
The option's value $V(F,\alpha)$ changes by $V_F\,dF+V_\alpha\,d\alpha$ to first order. The regression of
$d\alpha$ on $dF$ has slope $\rho\nu/F^\beta$, so the position $\Delta=V_F+V_\alpha\rho\nu/F^\beta$ removes
the part of $d\alpha$ correlated with $dF$ and leaves only the uncorrelated part, which no position in the
underlying can hedge.
\end{proposition}
\begin{proof}
$d\alpha=\nu\alpha\,dW^2=\nu\alpha(\rho\,dW^1+\sqrt{1-\rho^2}\,dW^\perp)$ and $dF=\alpha F^\beta dW^1$, so
$d\alpha=\frac{\rho\nu}{F^\beta}dF+\nu\alpha\sqrt{1-\rho^2}\,dW^\perp$. Substituting, the P\&L of
$V-\Delta F$ is $(V_F+V_\alpha\rho\nu/F^\beta-\Delta)\,dF+V_\alpha\nu\alpha\sqrt{1-\rho^2}\,dW^\perp$; its
variance is minimal when the first bracket vanishes. With $V=C_{\mathrm{BS}}(F,\sigma_B(F,\alpha))$,
$V_F=\Delta_{\mathrm{BS}}+\mathcal V\partial_F\sigma_B$ and $V_\alpha=\mathcal V\partial_\alpha\sigma_B$.
\end{proof}

On the example of the chapter, the one-year at-the-money call has a Black delta of 0.540, a Hagan delta of 0.580
and a \omterm{def:dv:sabr-and-smile-dynamics:mvdelta}{minimum-variance delta} of 0.461 (\cref{fig:dv:sabr-and-smile-dynamics:deltas}). With a negative correlation,
a rise in the forward tends to bring a fall in $\alpha$ and so in the call's value: the right hedge is smaller than
Black's. Hagan's delta moves the other way because it ignores the correlation and sees only the \omterm{def:dv:sabr-and-smile-dynamics:backbone}{backbone}.

\begin{omfigure}
\begin{tikzpicture}
\begin{axis}[width=10.5cm,height=4.8cm,
  xlabel={strike (\%)},ylabel={delta of a 1-year call},
  tick label style={font=\small},label style={font=\small},
  xmin=2,xmax=4.5,ymin=0,ymax=1,grid=major,grid style={gray!25},
  legend columns=3,legend style={font=\footnotesize,at={(0.5,-0.3)},anchor=north,draw=none,fill=none,/tikz/every even column/.append style={column sep=8pt}},
  table/col sep=comma]
\addplot[gray,thick,dashed] table[x=k,y=black]{figdata/derivatives/11-sabr-and-smile-dynamics/deltas.csv};
\addplot[omProp,thick] table[x=k,y=hagan]{figdata/derivatives/11-sabr-and-smile-dynamics/deltas.csv};
\addplot[omThm,very thick] table[x=k,y=bartlett]{figdata/derivatives/11-sabr-and-smile-dynamics/deltas.csv};
\legend{Black,Hagan,minimum-variance}
\end{axis}
\end{tikzpicture}

\omcaption{Three deltas of one-year calls on a 3\% forward by strike ($\beta=0.5$, $\rho=-0.6$, $\nu=0.5$).
The \omterm{def:dv:sabr-and-smile-dynamics:mvdelta}{minimum-variance delta} is the smallest, because a rise in the forward is expected to come with a fall in
volatility. Data: the tutorial.}\label{fig:dv:sabr-and-smile-dynamics:deltas}
\end{omfigure}

How much does it matter? The tutorial simulates one day of SABR dynamics on 200\,000 paths and reprices the
call with the formula at the new forward and $\alpha$. The standard deviation of the hedged one-day P\&L is
$6.67\times10^{-5}$ with the Black delta, $7.46\times10^{-5}$ with Hagan's and $5.99\times10^{-5}$ with the
\omterm{def:dv:sabr-and-smile-dynamics:mvdelta}{minimum-variance delta}: 19\% less variance than Black, and 25\% more with Hagan's. What remains is the part of the
volatility's move that is independent of the forward, hedged only with \omterm{def:dv:greeks-and-the-hedging-pnl:greeks}{vega} (another option).

\section{Limits: wings and long expiries}

Hagan's formula is an expansion in $T$ and in $\ln(F/K)$. It is accurate for the expiries and strikes of the
liquid swaption market, and it fails in two known ways. At long expiries with a large \omterm{def:dv:stochastic-volatility:sv}{volatility of volatility}
it produces call prices that are not convex at low strikes: a negative density
(\cref{fig:dv:sabr-and-smile-dynamics:density}). And its zero-order term was shown to be slightly wrong
(Obłój's correction), which matters in the wings. The fixes used in practice are to solve an effective
one-dimensional forward equation for the density with the same asymptotic accuracy (``arbitrage-free SABR''), or
to shift the forward (One Quant Book 6, chapter 5), or to price the wings with a different extrapolation.

\begin{omfigure}
\begin{tikzpicture}
\begin{axis}[width=10.5cm,height=4.8cm,
  xlabel={strike (\%)},ylabel={density from Hagan's prices},
  tick label style={font=\small},label style={font=\small},
  xmin=0,xmax=3,ymin=-130,ymax=60,grid=major,grid style={gray!25},
  table/col sep=comma]
\addplot[omThm,very thick] table[x=k,y=density]{figdata/derivatives/11-sabr-and-smile-dynamics/density.csv};
\draw[gray,thick] (axis cs:0,0) -- (axis cs:3,0);
\node[font=\footnotesize,anchor=west,align=left] at (axis cs:1.7,-80) {negative density:\\a butterfly arbitrage};
\end{axis}
\end{tikzpicture}

\omcaption{The Breeden--Litzenberger density implied by Hagan's formula for a ten-year option on a 3\% forward
($\beta=0.5$, $\rho=-0.3$, $\nu=0.6$, 20\% at the money). Below about 1.6\% it is negative: the formula prices
butterflies there below zero. Data: the tutorial.}\label{fig:dv:sabr-and-smile-dynamics:density}
\end{omfigure}

\section{Tutorial: SABR from formula to hedge}\label{tut:dv:sabr-and-smile-dynamics:tut}

\begin{tutorial}
\textbf{Goal.} Implement Hagan's formula, calibrate SABR with $\alpha$ tied to the at-the-money volatility, compare
three deltas, and measure their hedging errors by simulation. \textbf{End state:} the four figures of the chapter
and the numbers of the weekend problem.
\begin{tutsteps}
\item \textbf{The formula}, with the at-the-money limit of $z/x(z)$ handled explicitly:
\omcode{code/firm/sabr/firm_sabr.py}{23}{37}{Hagan's Black-volatility formula.}\label{lst:dv:sabr-and-smile-dynamics:hagan}
\item \textbf{The three deltas}, the correlation correction coming from $\partial\sigma_B/\partial\alpha$:
\omcode{code/firm/sabr/firm_sabr.py}{76}{89}{Black, Hagan and minimum-variance deltas.}\label{lst:dv:sabr-and-smile-dynamics:deltas}
\item \textbf{Run} \texttt{dv\_sabr.hedge\_experiment()}, \texttt{equity\_fit()}, \texttt{wing\_density()} and
\texttt{fig\_sabr.py}.
\end{tutsteps}
\textbf{What to change next.} Set $\rho=+0.3$ and compare the three deltas' hedging errors (the \omterm{def:dv:sabr-and-smile-dynamics:mvdelta}{minimum-variance
delta} then nearly equals Black's); then fit the equity smile with $\beta=0.5$ and compare $\rho$ and $\nu$ with the
$\beta=1$ fit.
\end{tutorial}

\section{Build: the SABR smile}\label{bld:dv:sabr-and-smile-dynamics:sabr}

\begin{build}
\bfield{Purpose} The miniature firm's smile for rates and currency options, and the parametrisation that One Quant
Book 6 extends to the volatility cube.
\bfield{Interface} \texttt{hagan\_lognormal(f, k, t, alpha, beta, rho, nu)}; \texttt{alpha\_from\_atm(f, t,
sigma\_atm, beta, rho, nu)}; \texttt{calibrate(f, t, strikes, vols, beta, atm\_vol) -> ((alpha, rho, nu), rmse)};
\texttt{normal\_vol}; \texttt{deltas(\ldots) -> black, hagan, bartlett}; \texttt{one\_day\_hedge\_errors}.
\bfield{Rules} $\beta$ is an input, not a fitted parameter; $\alpha$ is always solved from the at-the-money volatility
(the smallest positive root of the cubic); report the density of the wings, and do not use the formula where it is
negative.
\bfield{Acceptance tests} \texttt{code/firm/sabr/tests/}: at $\nu\to0$ and $\beta=1$ the smile is flat at $\alpha$; the
at-the-money volatility is reproduced exactly; the fit recovers known parameters; the \omterm{def:dv:sabr-and-smile-dynamics:mvdelta}{minimum-variance delta} has the
lowest simulated hedging error for negative correlations.
\bfield{Stretch} The arbitrage-free density of Hagan and his coauthors (2014), solved on a grid, and Obłój's corrected
leading term.
\end{build}

\begin{omsources}
\item P.~S.~Hagan, D.~Kumar, A.~S.~Lesniewski and D.~E.~Woodward, ``Managing smile risk'', \emph{Wilmott} (2002).
\item B.~Bartlett, ``Hedging under SABR model'', \emph{Wilmott} (July/August 2006) 2--4.
\item J.~Obłój, ``Fine-tune your smile: correction to Hagan et al.'', \emph{Wilmott} (May 2008) 102--109.
\item P.~S.~Hagan, D.~Kumar, A.~S.~Lesniewski and D.~E.~Woodward, ``Arbitrage-free SABR'', \emph{Wilmott} 69 (2014)
60--75.
\end{omsources}

\section{Exercises}

\begin{exercise}[$\star$]\label{exo:dv:sabr-and-smile-dynamics:1}
With $\beta=1$ and $\nu=0$, what is the SABR smile? With $\beta=0$ and $\nu=0$?
\end{exercise}

\begin{exercise}[$\star$]\label{exo:dv:sabr-and-smile-dynamics:2}
Using the leading term $\sigma_{\mathrm{ATM}}\approx\alpha F^{\beta-1}$ with $\alpha$ set for 20\% at a 3\% forward, give
the at-the-money volatility at a 2.5\% forward for $\beta=0$, 0.5 and 1.
\end{exercise}

\begin{exercise}[$\star$]\label{exo:dv:sabr-and-smile-dynamics:3}
Why does SABR not need mean reversion when it is calibrated expiry by expiry?
\end{exercise}

\begin{exercise}[$\star\star$]\label{exo:dv:sabr-and-smile-dynamics:4}
Derive the correction term of the \omterm{def:dv:sabr-and-smile-dynamics:mvdelta}{minimum-variance delta} from the regression of $d\alpha$ on $dF$.
\end{exercise}

\begin{exercise}[$\star\star$]\label{exo:dv:sabr-and-smile-dynamics:5}
$\beta$ and $\rho$ both produce skew. How would you choose $\beta$ for a currency pair, and what does the wrong choice
cost?
\end{exercise}

\begin{exercise}[$\star\star$]\label{exo:dv:sabr-and-smile-dynamics:6}
Give the normal volatility of the 2\%, 3\% and 4\% strikes in \cref{ex:dv:sabr-and-smile-dynamics:rates} and explain why
the normal smile is flatter than the Black smile.
\end{exercise}

\begin{exercise}[$\star\star\star$]\label{exo:dv:sabr-and-smile-dynamics:7}
\emph{Coding.} Repeat the one-day hedging experiment with $\rho=+0.3$, $\nu=0.4$: which delta is best, and by how much?
\end{exercise}

\begin{exercise}[$\star\star\star$]\label{exo:dv:sabr-and-smile-dynamics:8}
\emph{Find the flaw.} ``Our SABR fit is perfect at every expiry, so its ten-year 1\% strike price is reliable.''
\end{exercise}

\section{Problem: Which Delta}

\begin{problem}[{Weekend problem --- the delta a smile model should use}]\label{pb:dv:sabr-and-smile-dynamics:1}
A rates desk is long one-year at-the-money calls on a 3\% forward (payer swaptions in normalised units), priced with
SABR: $\beta=0.5$, $\rho=-0.6$, $\nu=0.5$, 20\% at-the-money Black volatility, zero discounting. It must choose the
delta it hedges with.

\medskip\noindent\textbf{Part I --- The model.}
\begin{enumerate}
\item Solve for $\alpha$.
\item Give the Black volatility at the 2\%, 3\% and 4\% strikes.
\item Give the at-the-money normal volatility in basis points.
\item What does the \omterm{def:dv:sabr-and-smile-dynamics:backbone}{backbone} predict for the at-the-money volatility if the forward rises to 3.5\%?
\item How does the smile move when the forward rises?
\end{enumerate}

\medskip\noindent\textbf{Part II --- Three deltas.}
\begin{enumerate}[resume]
\item Give the Black, Hagan and \omterm{def:dv:sabr-and-smile-dynamics:mvdelta}{minimum-variance deltas} of the at-the-money call.
\item Why is Hagan's delta larger than Black's?
\item Why is the \omterm{def:dv:sabr-and-smile-dynamics:mvdelta}{minimum-variance delta} smaller than both?
\item Give the regression slope of $d\alpha$ on $dF$.
\item What risk remains after the minimum-variance hedge, and how is it hedged?
\end{enumerate}

\medskip\noindent\textbf{Part III --- The experiment.}
\begin{enumerate}[resume]
\item Give the standard deviations of the one-day hedged P\&L with each delta.
\item Give the variance of each relative to Black's.
\item Why does Hagan's delta do worse than Black's here?
\item What would the result be with $\rho=0$?
\item If daily hedging errors are independent, by how much does a 19\% lower daily variance reduce the
standard deviation of a year's \omterm{def:dv:greeks-and-the-hedging-pnl:hpnl}{hedging P\&L}?
\end{enumerate}

\medskip\noindent\textbf{Part IV --- Judgement.}
\begin{enumerate}[resume]
\item Which delta would you report to risk, and which would you trade?
\item How does the choice of $\beta$ change the deltas?
\item What if the market's smile dynamics differ from SABR's?
\item State the \emph{named result}: the reduction in the variance of the one-day hedging error from the
\omterm{def:dv:sabr-and-smile-dynamics:mvdelta}{minimum-variance delta} relative to the Black delta.
\item In one sentence: what does a smile model's delta depend on that its prices do not?
\end{enumerate}
\end{problem}

\section{Interview questions}

\begin{interviewq}[$\star$ \iqroles{trader, researcher}]\label{iq:dv:sabr-and-smile-dynamics:1}
Write the \omterm{def:dv:sabr-and-smile-dynamics:sabr}{SABR model}. What does each parameter do?
\end{interviewq}

\begin{interviewq}[$\star$ \iqroles{trader}]\label{iq:dv:sabr-and-smile-dynamics:2}
What is the \omterm{def:dv:sabr-and-smile-dynamics:backbone}{backbone}, and how would you estimate $\beta$ from market data?
\end{interviewq}

\begin{interviewq}[$\star\star$ \iqroles{researcher}]\label{iq:dv:sabr-and-smile-dynamics:3}
Derive the \omterm{def:dv:sabr-and-smile-dynamics:mvdelta}{minimum-variance delta} in a \omterm{def:dv:stochastic-volatility:sv}{stochastic volatility model}.
\end{interviewq}

\begin{interviewq}[$\star\star$ \iqroles{researcher, risk}]\label{iq:dv:sabr-and-smile-dynamics:4}
Why did \omterm{def:dv:local-volatility:model}{local volatility} give bad hedges in the rates market, according to the authors of SABR?
\end{interviewq}

\begin{interviewq}[$\star\star$ \iqroles{developer}]\label{iq:dv:sabr-and-smile-dynamics:5}
Your SABR implementation returns negative butterfly prices for long-dated low strikes. Why, and what do you do?
\end{interviewq}

\begin{interviewq}[$\star\star\star$ \iqroles{trader, researcher}]\label{iq:dv:sabr-and-smile-dynamics:6}
$\beta$ and $\rho$ fit the same smile equally well. Which trades distinguish them, and how would you hedge a book in
which they are uncertain?
\end{interviewq}
